\begin{proposition}[Vanishing of the normalized radial-loop defect]\hypertarget{tag-0312}{}\label{push:src:B.4.32}
\label{push:B432}
\StatementRecord{push:src:B.4.32}{proved}{0312}{\SeriesRef{push:src:B.4.31}, \SeriesRef{push:src:B.4.19}, \SeriesRef{push:src:B.4.4}, \SeriesRef{push:src:B.1.2}, \SeriesRef{push:src:B.4.30}; \citet[Theorem 6.3]{IKN2005E}}

Suppose a proper complex-analytic filling of a finite parameter
cover of the six-sphere family exists. Make the
modifications in \ref{push:src:B.4.31} and delete its finite exceptional set
from the marked base. Use the ordinary local origins of
\ref{push:src:B.4.19} as reference. There is a well-defined normalized
radial-loop defect of a local extending cap section,
\begin{equation}
 \lambda\in
 V:=\bigl(\Lambda\otimes\Q\bigr)/\Q\widehat\delta .
\label{pub:I:eq:B-4-32-1}
\end{equation}
It is independent of the section, of the finite parameter root
needed to obtain that section, and of the modifications used.
It patches to a flat section of the local system $V$
over the remaining ordinary base, and this section is zero.

In particular, let $t$ be the trace of an extending finite
multisection of degree $d$, and let $o$ be the reference ordinary
origin. The radial-loop defect of $t-do$ in $V$ is zero.
The same holds for an origin obtained from multiplied such
traces by bounded logarithmic differences along $r=0$.

\end{proposition}



\begin{proof}
\begin{enumerate}[label=\arabic*.,leftmargin=*]
\item Use the smooth marked model of Theorem~\ref{night:P2.4} and remove the finitely many base values which are images of vertical central components or critical values of the maps on the central branch strata. Every central component above the remaining ordinary base dominates the original $\mathbb P^1$. Its positive normal meridian has zero lattice class in $\Lambda/\mathbb Z\widehat\delta$ by Theorem~\ref{night:P2.6}.

\item The ordinary origin is exactly the one used there. Indeed, write $k=\log K/(2\pi i)$ with $K=t^4/(t-1)^3$ and $v=\log r/(2\pi i)$. The parameter descent of Theorem~\ref{night:P2.1} is $H=-k$, so subtracting the local origin $vHe$, with $e=\Pi_c\widehat\delta$ gives the ordinary gauge $\widetilde\zeta=\zeta+kv e$ of Proposition~\ref{push:src:B.4.19}. On an overlap, $k$ changes by an integer, and the resulting radial label is an integral multiple of $\widehat\delta$. The finite-root cancellation and the cusp comparison in that proposition have the same property. Thus the splitting used in Theorem~\ref{night:P2.6} computes precisely the radial-loop defect relative to the stated ordinary origins.

\item The connecting map is the endpoint period of a continued additive lift. With $T_r=I-mE$, $E=\widehat\delta\otimes\gamma$ and $E^2=0$, changing the lift adds $(T_r-I)\Lambda\subset\mathbb Z\widehat\delta$, so its image $\nu_r$ in $V$ is well defined. Under $r=R^e$ its label changes by
\[
 \sum_{j=0}^{e-1}T_r^jv
 =ev-\frac{me(e-1)}2Ev
 \equiv ev\pmod{\mathbb Q\widehat\delta}.
\]
Consequently the normalized defect of a section $b$ on that root cover is $e^{-1}\nu_R(b-o(R^e))$.

\item Lift an extending section to a common proper resolution. For a fixed ordinary base point the lift extends over the parameter origin: normalize the dominating curve in its proper graph, which has degree one over the disc. At its limit choose a normal-crossing chart in which
\[
 r=\text{unit}\prod_{i=1}^k x_i^{a_i},\qquad a_i>0.
\]
On the section, $x_i(R)=R^{b_i}u_i(R)$ with $u_i$ a holomorphic unit and $b_i\ge0$. Thus $e=\sum_i a_i b_i$, and its positive $R$-loop is the sum of $b_i$ positive component meridians. The unit loops contract in the chart. In the quotient presentation of Theorem~\ref{night:P2.6}, the lattice label is the same sum of the component labels, each zero modulo $\widehat\delta$. Hence the normalized radial defect is zero. This argument applies to every extending section and every root, and is preserved under common proper refinements. In a varying ordinary base chart the labels are locally constant, so the vanishing also holds across the isolated points where a simultaneous section lift has indeterminacy.

\item The connecting homomorphism is additive. Split a finite multisection after a cyclic root at a generic ordinary point, add the zero normalized defects of its branches, and use the preceding root formula. Its trace of degree $d$ therefore satisfies $\nu_r(t-do)=0$ in $V$. These zero local classes patch under the lattice transport, independently of the section and all choices. Finally, by the bounded-character presentation of Proposition~\ref{push:src:B.4.27}, a bounded logarithmic difference along the strict parameter axis has characters in $\mathcal O^*\times\mathcal M_r^{\mathrm{gp}}$: its first character has zero winding and its second contributes only a multiple of $\widehat\delta$. Adding it to a multiplied trace preserves the zero class. This proves every stated assertion.
\end{enumerate}
\end{proof}